\documentclass{beamer}

\def\Tiny{\fontsize{6pt}{6pt}\selectfont}
\def\supertiny{\fontsize{4pt}{4pt}\selectfont}

\mode<presentation>
{
  \usetheme{Warsaw}
  % \setbeamercovered{transparent}
  \usecolortheme{crane}
}


\usepackage{graphicx, ifthen, listings,fancyvrb}

\usepackage[czech]{babel}
% \usefonttheme{professionalfonts}
\usepackage{times}
\usepackage{amsmath}
\usepackage[utf8]{inputenc}
\usepackage{wrapfig}

\usepackage[T1]{fontenc}

\lstset{ basicstyle=\tiny, stringstyle=\ttfamily, showstringspaces=false }

\everymath{\displaystyle}

\setbeamerfont{frametitle}{size=\large}
\setbeamerfont{subsection in toc}{size=\scriptsize}

\makeatletter\newenvironment{blackbox}{%
   \begin{lrbox}{\@tempboxa}\begin{minipage}{0.95\columnwidth}}{\end{minipage}\end{lrbox}%
   \colorbox{black}{\usebox{\@tempboxa}}
}\makeatother

\title[IMF (4)]{Informatika pro moderní fyziky (4)\\práce s tabulkami, hromadné zpracování dat, rozšíření RubyGems, ladění programů}

\author[Franti\v{s}ek HAVL\r{U}J, ORF ÚJV Řež]{Franti\v{s}ek HAVL\r{U}J\\{\scriptsize \emph{e-mail: haf@ujv.cz}}}

\date{akademický rok 2022/2023, 15. října 2026}

\institute[ORF ÚJV Řež]
{ÚJV Řež\\oddělení Reaktorové fyziky a podpory palivového cyklu}

\AtBeginSection[]
{
\begin{frame}<beamer>
\frametitle{Obsah}
\tableofcontents[currentsection,hideothersubsections]
\end{frame}
}

\begin{document}

\begin{frame}
  \titlepage
\end{frame}

\begin{frame}
  \tableofcontents
\end{frame}

\section{Co jsme se naučili minule}

\begin{frame}{}
  \begin{itemize}
    \item zpracování většího množství dat
    \item
  \end{itemize}
\end{frame}

\section{Rozšíření Ruby: RubyGems a Bundler}

\begin{frame}{Knihovny (gemy) jsou základ}
  \begin{itemize}
    \item existují mnohá rozšíření, tzv. knihovny -- v ruby se jim říká \textbf{rubygems}
    \item aktuálně nás zajímá něco na práci s excelovskými soubory
    \item gemy jdou sice instalovat na systémové úrovni, ale z toho je pak zase jenom neštěstí
    \item použijeme radši \textbf{bundler}, správce gemů pro každého: vyřeší za nás závislosti a postará se o snadnou instalaci
  \end{itemize}
\end{frame}

\begin{frame}{Máme bundler?}
  \begin{itemize}
    \item otestujeme rubygems: \texttt{gem -v}
    \item pokud není, zapláčeme, protože jsme asi špatně nainstalovali Ruby
    \item otestujeme bundler: \texttt{bundle -v}
    \item pokud bundler není, doinstalujeme \texttt{gem install bundler}
  \end{itemize}
\end{frame}

\begin{frame}{Jak na to}
  \begin{itemize}
    \item najdu si, která knihovna mě zajímá (třeba na rubygems.org nebo kdekoli jinde): my bychom rádi \textbf{rubyXL} \texttt{https://github.com/weshatheleopard/rubyXL}
    \item vytvořím si prázdný \texttt{Gemfile} -- tam se specifikuje, které gemy chci používat: \texttt{bundle init}
    \item do gemfilu -- je to normální Ruby skript! -- dopíšu\\ \texttt{gem "rubyXL"}
    \item nainstaluju: \texttt{bundle install}
  \end{itemize}
\end{frame}

\begin{frame}{Jak použít?}
  \begin{itemize}
    \item na začátku svého skriptu pak musím nahrát bundler:
    \item \texttt{require "bundler/setup"}
    \item a teď už můžu nahrát jakýkoli gem:
    \item \texttt{require "rubyXL"}
  \end{itemize}
\end{frame}

\subsection{Vytvoříme excelovskou tabulku}

\begin{frame}[fragile]{RTFM, RTFM, RTFM}
  \begin{itemize}
    \item na stránkách \texttt{rubyXL} se nachází spousta příkladů a návodů -- https://github.com/weshatheleopard/rubyXL
    \item kromě toho má i slušnou dokumentaci (GIYF / "rubyxl docs") -- http://www.rubydoc.info/gems/rubyXL/3.3.15
    \item napoprvé navedu do začátku:
  \end{itemize}
  {\scriptsize
  \begin{verbatim}
    workbook = RubyXL::Workbook.new
    worksheet = workbook[0]
    worksheet.add_cell(0, 0, "A1")
    workbook.write("data.xlsx")
  \end{verbatim}
  }
\end{frame}

\begin{frame}[fragile]{Kompletní test}
  ... se nachází v data/bundler.
  \begin{block}{Gemfile}
  {\scriptsize
  \begin{verbatim}
source "https://rubygems.org"
gem "rubyXL"
  \end{verbatim}
  }
  \end{block}
  \begin{block}{Ruby skript}
  {\scriptsize
  \begin{verbatim}
require "bundler/setup"
require "rubyXL"

workbook = RubyXL::Workbook.new
worksheet = workbook[0]
worksheet.sheet_name = "Data test"
worksheet.add_cell(0, 0, "A1")
workbook.write("data.xlsx")
  \end{verbatim}
  }
  \end{block}
\end{frame}

\begin{frame}{Jednoduché cvičení}
  \begin{itemize}
    \item použijte soubor \texttt{numbers/data\_two\_1.csv}
    \item vytvořte excelovský soubor se dvěma listy, na obou bude sloupec 1, sloupec 2 a součet
    \item na jednom součet bude jako číslo (sečte to váš skript)
    \item na druhém bude součet jako excelovský vzorec
  \end{itemize}
\end{frame}


\section{Kde je chyba?}

\begin{frame}{Ladění programů}
  \begin{itemize}
    \item v každém programu je aspoň jedna chyba
    \item není důležité nedělat chyby, ale je nutné je umět najít
    \item když si program/Ruby na něco stěžuje, tak si to přečtěte, jinak se nic nedozvíte
    \item pokud nepoznám, v čem je chyba, jsem bezezbytku ztracen
    \item následují tři úlohy, kde je úkolem najít všechny chyby
    \item z didaktických důvodů postupujte metodou tupého spouštění a postupného opravování
  \end{itemize}
\end{frame}

\section{Zpracování textu}

\subsection{Obecný rozbor}

\begin{frame}{Problém č. 2: mnoho výpočtů, inženýrova smrt}
  \begin{block}{Zadání}
    Při přípravě základního kritického experimentu je pomocí MCNP potřeba najít kritickou polohu regulační tyče R2.

    Jak se tato poloha změní při změně polohy tyče R1?
  \end{block}
\end{frame}

\begin{frame}{Co máme k~dispozici?}
  \begin{block}{MCNP}
    Pokud připravíme vstupní soubor (v~netriviální formě obsahující polohy regulačních tyčí R1 a R2), spočítá nám keff.
  \end{block}
  Potřebovali bychom ale něco na:
  \begin{enumerate}
    \item vytvoření velkého množství vstupních souborů
    \item extrakci keff z~výstupních souborů
    \item popřípadě na vyhodnocení získaných poloh tyčí a keff
  \end{enumerate}
\end{frame}

\begin{frame}{Pracovní postup}
  \begin{enumerate}
    \item načíst keff z~výstupního souboru MCNP
    \pause
    \item vygenerovat potřebné vstupní soubory
    \pause
    \item vyrobit BAT soubor na spuštění výpočtů
    \pause
    \item načíst výsledky ze všech výstupních souborů do jedné tabulky
  \end{enumerate}
\end{frame}

\subsection{Načítání výstupního souboru}

\begin{frame}[fragile]{}
  Nejprve najdeme, kde je ve výstupu z~MCNP žádané keff:
  {\tiny
  \begin{verbatim}
.....
         the k(trk length) cycle values appear normally distributed at the 95 percent confidence level

-----------------------------------------------------------------------------------------------------------------------------------
|                                                                                                                                 |
| the final estimated combined collision/absorption/track-length keff = 1.00353 with an estimated standard deviation of 0.00024   |
|                                                                                                                                 |
| the estimated 68, 95, & 99 percent keff confidence intervals are 1.00329 to 1.00377, 1.00305 to 1.00400, and 1.00289 to 1.00416 |
|                                                                                                                                 |
| the final combined (col/abs/tl) prompt removal lifetime = 1.0017E-04 seconds with an estimated standard deviation of 2.7364E-08 |
.....
  \end{verbatim}
  }
\end{frame}

\begin{frame}{Algoritmus}
  \begin{enumerate}
    \item najít řádek s~keff
    \pause
    \item vytáhnout z~něj keff, takže například:
    \pause
    \item rozdělit podle rovnítka
    \pause
    \item druhou část rozdělit podle mezer
    \pause
    \item vzít první prvek
  \end{enumerate}
\end{frame}


\begin{frame}[fragile]{Realizace (1/5)}
  \scriptsize
  \begin{verbatim}
    keff = nil

    File.foreach("c1_1o") do |line|

    end

    puts keff
  \end{verbatim}
\end{frame}

\begin{frame}[fragile]{Realizace (2/5)}
  \scriptsize
  \begin{verbatim}
    keff = nil

    File.foreach("c1_1o") do |line|
      if line.include?("final estimated combined")

      end
    end

    puts keff
  \end{verbatim}
\end{frame}

\begin{frame}[fragile]{Realizace (3/5)}
  \scriptsize
  \begin{verbatim}
    keff = nil

    File.foreach("c1_1o") do |line|
      if line.include?("final estimated combined")
        a = line.split("=")

      end
    end

    puts keff
  \end{verbatim}
\end{frame}

\begin{frame}[fragile]{Realizace (4/5)}
  \scriptsize
  \begin{verbatim}
    keff = nil

    File.foreach("c1_1o") do |line|
      if line.include?("final estimated combined")
        a = line.split("=")
        b = a[1].split

      end
    end

    puts keff
  \end{verbatim}
\end{frame}

\begin{frame}[fragile]{Realizace (5/5)}
  \scriptsize
  \begin{verbatim}
    keff = nil

    File.foreach("c1_1o") do |line|
      if line.include?("final estimated combined")
        a = line.split("=")
        b = a[1].split
        keff = b[0]
      end
    end

    puts keff
  \end{verbatim}
\end{frame}


\section{Úvod do OOP}

\begin{frame}{Co je OOP (1)}
  \begin{itemize}
    \item zatím jsme používali tzv. procedurální programování – máme data a pak procedury/funkce
    \item Ruby je nicméně čistě objektový jazyk, i když jsme se tomu zatím spíš vyhýbali
    \item objekt je entita, která má \emph{vlastnosti} (properties) a \emph{metody} (methods) a poskytuje okolnímu světu nějaké \emph{rozhraní} (interface)
  \end{itemize}
\end{frame}

\begin{frame}{Co je OOP (2)}
  \begin{itemize}
    \item většinou se jako hlavní důvody pro OOP uvádí polymorfismus a dědičnost (inheritance), ale to hlavní je posun uvažování od dat a operací nad nimi k ``inteligentním'' objektům – a spolu s tím zapouzdření (encapsulation)
    \item je potřeba si uvědomit, že všechna `zázračná' paradigmata v programování jsou pouze odlišné formalismy, může nám to hodně pomoct a stojí za to se tomu věnovat, na druhou stranu je potřeba se z toho nezbláznit a nedělat z toho náboženství
  \end{itemize}
\end{frame}

\begin{frame}{Třída}
  \begin{itemize}
    \item `typ' objektu - definuju, jak se objekty chovají
    \item neboli definuju metody
    \item zjednodušeně řečeno jsou to jen metody a ne data
    \item potkali jsme třeba \texttt{File}, \texttt{CSV} (a aniž bychom to věděli, tak \texttt{Fixnum}, \texttt{Array}, \texttt{Hash})
  \end{itemize}
\end{frame}

\begin{frame}[containsverbatim]{Instance}
  \begin{itemize}
    \item konkrétní objekt – má svoje data
    \item většinou vytváříme pomocí \texttt{Class.new}, resp \texttt{Class.new(arg1, arg2, ...)}
    \item k datům přistupuju pomocí \emph{instance variables} -- \verb|@data| -- ta nikdo jiný zvenku nevidí
    \item můžu definovat \texttt{attr\_reader}, \texttt{attr\_writer}, \texttt{attr\_accessor} pro přístup k instance variables zvenčí
    \item metody se definují podobně jako normální funkce a volají se obvyklým způsobem \texttt{obj.method}
    \item leckdy se hodí speciální metoda \texttt{initialize}, tzv. \emph{konstruktor} -- volá se při \texttt{new} a nastavují se zde výchozí hodnoty vlastností
  \end{itemize}
\end{frame}

\begin{frame}{Class method}
  \begin{itemize}
    \item můžu mít i metody, které nepatří k žádné instanci – nepracují s žádnými daty
    \item dávat je do třídy má pak smysl pouze organizační, nemá to pak žádnou reálnou výhodu proti obyčejným funkcím
    \item např \texttt{File.foreach} nebo \texttt{CSV.read}
  \end{itemize}
\end{frame}

\begin{frame}[containsverbatim]{Jak to vypadá v Ruby}
  \begin{Verbatim}[fontsize=\tiny]
  class Table
    def initialize
      @data = {}
    end
    def get(key)
      @data[key]
    end
    def set(key, value)
      @data[key] = value
    end
    def print
      @data.each do |key, value|
        puts "#{key} #{value}"
      end
    end
  end

  t = Table.new
  t.set("a", 123)
  t.print
  \end{Verbatim}
\end{frame}

\begin{frame}[containsverbatim]{Úkol: implementujte třídu MCNPFile z příkladu}
  \begin{verbatim}
    f = MCNPOutput.new("c1_1o")
    puts f.keff
    puts f.keff_uncertainty
  \end{verbatim}
\end{frame}

\section{Dáme věci dohromady}

\subsection{Zápis všech výsledků do tabulky}

\begin{frame}{Jak na to}
  Máme všechno, co potřebujeme:
  \begin{itemize}
    \item načtení keff z~jednoho výstupního souboru (třída \texttt{MCNPOutput})
    \item procházení adresáře (\texttt{Dir.each})
    \item zápis do souboru (\texttt{File.open} s~parametrem \texttt{w} anebo \texttt{File.write})
  \end{itemize}
  Takže už to stačí jen vhodným způsobem spojit dohromady!
\end{frame}

\begin{frame}[fragile]{Realizace}
  \scriptsize
  \begin{verbatim}
    Dir["*o"].each do |filename|
      f = MCNPOutput.new(filename)
      puts format("%-30s %10.5f", filename, f.keff)
    end
  \end{verbatim}
\end{frame}

\begin{frame}[fragile]{Výstup}
  Výsledkem je perfektní tabulka:
  {\scriptsize
  \begin{verbatim}
    outputs/c_0_0o      0.94800
    outputs/c_0_10o     0.99800
    outputs/c_0_1o      0.94850
    outputs/c_0_2o      0.95000
    outputs/c_0_3o      0.95250
    outputs/c_0_4o      0.95600
    ...
  \end{verbatim}}
  Hloupé je, že nikde nemáme tu polohu tyčí.
\end{frame}


\begin{frame}[fragile]{Chytrá horákyně}
  ... by jistě vyrobila toto:
  {\scriptsize
  \begin{verbatim}
    0  0    0.94800
    0 10    0.99800
    0  1    0.94850
    0  2    0.95000
    0  3    0.95250
    0  4    0.95600
    ...
  \end{verbatim}
  }
  Nápovědou je funkce \texttt{split} (podle podtržítka) a metoda \texttt{to\_i}
\end{frame}

\begin{frame}[fragile]{Realizace chytré horákyně}
  \scriptsize
  \begin{verbatim}
    Dir["outputs/*o"].each do |filename|
      keff = nil

      File.foreach(filename) do |line|
        if line.include?("final estimated combined")
          a = line.split("=")
          b = a[1].split
          keff = b[0]
        end
      end

      s = filename.split("_")
      puts "#{s[1].to_i} #{s[2].to_i} #{keff}"
    end
  \end{verbatim}
\end{frame}

\begin{frame}{A to je vše, přátelé!}
  \begin{center}
    \includegraphics[width=0.8\textwidth]{looney_tunes}
  \end{center}
\end{frame}

\end{document}
